% LaTeX Template for AI Problem Analysis Output (v2.0)
% Strategic Analysis of AMC12B 2023 Problem 24
% IMPORTANT: Use LaTeX commands for all mathematical symbols, NOT unicode characters

\documentclass[12pt]{article}
\usepackage[utf8]{inputenc}
\usepackage{amsmath, amssymb, amsthm}
\usepackage{geometry}
\usepackage{xcolor}
\usepackage[most]{tcolorbox}
\usepackage{enumitem}
\usepackage{fancyhdr}

% Page setup
\geometry{margin=1in}
\setlength{\headheight}{14.5pt}
\pagestyle{fancy}
\fancyhf{}
\rhead{AMC12B 2023 Problem 24 Analysis}
\lhead{\today}

% Color definitions
\definecolor{problemconfig}{RGB}{230, 242, 255}    % Light blue
\definecolor{strategic}{RGB}{255, 230, 242}       % Light pink
\definecolor{tactical}{RGB}{242, 255, 230}        % Light green
\definecolor{techniques}{RGB}{255, 242, 230}      % Light yellow
\definecolor{verification}{RGB}{255, 230, 230}    % Light red
\definecolor{solution}{RGB}{230, 255, 255}        % Light cyan
\definecolor{competition}{RGB}{242, 242, 255}     % Light purple
\definecolor{proofwriting}{RGB}{240, 230, 255}    % Light lavender
\definecolor{gold}{RGB}{218, 165, 32}

% Box styles
\newtcolorbox{problemconfigbox}{
    colback=problemconfig,
    colframe=blue,
    title=Problem Configuration Analysis,
    fonttitle=\bfseries,
    arc=3pt,
    boxrule=1pt,
    breakable
}

\newtcolorbox{strategicbox}{
    colback=strategic,
    colframe=purple,
    title=Strategic Framework Application,
    fonttitle=\bfseries,
    arc=3pt,
    boxrule=1pt,
    breakable
}

\newtcolorbox{tacticalbox}{
    colback=tactical,
    colframe=green,
    title=Tactical Methodology Selection,
    fonttitle=\bfseries,
    arc=3pt,
    boxrule=1pt,
    breakable
}

\newtcolorbox{techniquesbox}{
    colback=techniques,
    colframe=orange,
    title=Conversion Techniques Application,
    fonttitle=\bfseries,
    arc=3pt,
    boxrule=1pt,
    breakable
}

\newtcolorbox{verificationbox}{
    colback=verification,
    colframe=red,
    title=Verification Strategy Design,
    fonttitle=\bfseries,
    arc=3pt,
    boxrule=1pt,
    breakable
}

\newtcolorbox{solutionbox}{
    colback=solution,
    colframe=cyan,
    title=Solution Roadmap,
    fonttitle=\bfseries,
    arc=3pt,
    boxrule=1pt,
    breakable
}

\newtcolorbox{competitionbox}{
    colback=competition,
    colframe=purple,
    title=Competition Strategy Integration,
    fonttitle=\bfseries,
    arc=3pt,
    boxrule=1pt,
    breakable
}

\newtcolorbox{proofbox}{
    colback=proofwriting,
    colframe=violet,
    title=Proof Structure and Justification (COMC Part C),
    fonttitle=\bfseries,
    arc=3pt,
    boxrule=1pt,
    breakable
}

\newtcolorbox{keyinsightbox}{
    colback=yellow!20,
    colframe=gold,
    title=Key Insight,
    fonttitle=\bfseries,
    arc=3pt,
    boxrule=2pt,
    leftrule=5pt,
    breakable
}

\newtcolorbox{warningbox}{
    colback=red!10,
    colframe=red!50,
    title=Warning: Common Pitfall,
    fonttitle=\bfseries,
    arc=3pt,
    boxrule=2pt,
    leftrule=5pt,
    breakable
}

\newtcolorbox{tipbox}{
    colback=green!10,
    colframe=green!50,
    title=Pro Tip,
    fonttitle=\bfseries,
    arc=3pt,
    boxrule=2pt,
    leftrule=5pt,
    breakable
}

\begin{document}

\title{AMC12B 2023 Problem 24\\Strategic Analysis}
\author{Competition Math Framework v2.1}
\date{\today}
\maketitle

\section*{Problem Statement}

Suppose that $a$, $b$, $c$ and $d$ are positive integers satisfying all of the following relations:

\begin{align*}
abcd &= 2^6 \cdot 3^9 \cdot 5^7\\
\text{lcm}(a,b) &= 2^3 \cdot 3^2 \cdot 5^3\\
\text{lcm}(a,c) &= 2^3 \cdot 3^3 \cdot 5^3\\
\text{lcm}(a,d) &= 2^3 \cdot 3^3 \cdot 5^3\\
\text{lcm}(b,c) &= 2^1 \cdot 3^3 \cdot 5^2\\
\text{lcm}(b,d) &= 2^2 \cdot 3^3 \cdot 5^2\\
\text{lcm}(c,d) &= 2^2 \cdot 3^3 \cdot 5^2
\end{align*}

What is $\gcd(a,b,c,d)$?

$$\textbf{(A)}~30\qquad\textbf{(B)}~45\qquad\textbf{(C)}~3\qquad\textbf{(D)}~15\qquad\textbf{(E)}~6$$

\begin{problemconfigbox}
\subsection*{A. Problem Classification}
\begin{itemize}[leftmargin=*]
    \item \textbf{Problem Type}: To Find (GCD of four numbers)
    \item \textbf{Competition Context}: AMC 12B 2023, Problem 24
    \item \textbf{Difficulty Band}: Late-band (P24) - Extremely challenging
    \item \textbf{Mathematical Domain}: Number Theory (GCD/LCM with Prime Factorization)
    \item \textbf{Estimated Time}: 7-10 minutes (requires systematic analysis of exponents)
\end{itemize}

\subsection*{B. Problem Structure Identification}
\begin{itemize}[leftmargin=*]
    \item \textbf{Given Information}: 
    \begin{itemize}
        \item Product of four numbers: $abcd = 2^6 \cdot 3^9 \cdot 5^7$
        \item Six LCM relations for pairs: $\text{lcm}(a,b)$, $\text{lcm}(a,c)$, etc.
        \item All expressed in prime factorization form
        \item Only three primes involved: 2, 3, 5
    \end{itemize}
    \item \textbf{Constraints}: 
    \begin{itemize}
        \item All numbers are positive integers
        \item Product constraint gives sum of exponents
        \item LCM constraints give max of exponents for pairs
        \item System must be consistent
    \end{itemize}
    \item \textbf{Objective}: 
    \begin{itemize}
        \item Find $\gcd(a,b,c,d)$
        \item GCD determined by minimum exponent of each prime
    \end{itemize}
    \item \textbf{Hidden Assumptions}: 
    \begin{itemize}
        \item Key insight: use exponent notation $\nu_p(x)$ = exponent of prime $p$ in $x$
        \item For each prime separately, analyze constraints
        \item $\text{lcm}(x,y)$ has exponent $\max\{\nu_p(x), \nu_p(y)\}$
        \item Product gives $\nu_p(a) + \nu_p(b) + \nu_p(c) + \nu_p(d)$
        \item GCD has exponent $\min\{\nu_p(a), \nu_p(b), \nu_p(c), \nu_p(d)\}$
    \end{itemize}
    \item \textbf{Answer Choice Leverage}: 
    \begin{itemize}
        \item (A) 30 = $2 \cdot 3 \cdot 5$: all three primes
        \item (B) 45 = $3^2 \cdot 5$: powers of 3 and 5
        \item (C) 3: only prime 3
        \item (D) 15 = $3 \cdot 5$: primes 3 and 5
        \item (E) 6 = $2 \cdot 3$: primes 2 and 3
    \end{itemize}
    \item \textbf{Proof Requirements}: Not a proof problem; systematic calculation
\end{itemize}

\subsection*{C. Strategic Orientation}
\begin{itemize}[leftmargin=*]
    \item \textbf{Hypothesis}: Analyze each prime separately using exponent notation
    \item \textbf{Conclusion}: Only prime 3 has minimum exponent $\geq 1$, so $\gcd = 3$
    \item \textbf{Notation}: 
    \begin{itemize}
        \item $\nu_p(x)$ = exponent of prime $p$ in prime factorization of $x$
        \item For $p \in \{2, 3, 5\}$, find $\nu_p(a), \nu_p(b), \nu_p(c), \nu_p(d)$
        \item GCD = $2^{\min\{\nu_2\}} \cdot 3^{\min\{\nu_3\}} \cdot 5^{\min\{\nu_5\}}$
    \end{itemize}
    \item \textbf{Intuitive Approach}: 
    \begin{itemize}
        \item For each prime, set up system of constraints
        \item Use max constraints from LCMs
        \item Use sum constraint from product
        \item Solve for individual exponents
        \item Find minimum exponent for each prime
    \end{itemize}
    \item \textbf{Cross-Domain Potential}: 
    \begin{itemize}
        \item Number theory: GCD/LCM properties
        \item Algebra: system of equations with max/min
        \item Logic: constraint satisfaction
    \end{itemize}
\end{itemize}
\end{problemconfigbox}

\begin{strategicbox}
\subsection*{A. Psychological Strategies}
\begin{itemize}[leftmargin=*]
    \item \textbf{Mental Toughness}: P24 is extremely difficult—stay systematic
    \item \textbf{Creativity Cultivation}: Exponent notation simplifies the problem
    \item \textbf{Iterative Refinement}: Solve for one prime at a time
\end{itemize}

\subsection*{B. Investigation Strategies}
\begin{itemize}[leftmargin=*]
    \item \textbf{Startup Orientation}: 
    \begin{itemize}
        \item Recognize all numbers in prime factorization form
        \item Use exponent notation to simplify
        \item Attack one prime at a time
    \end{itemize}
    \item \textbf{Get Your Hands Dirty}: 
    \begin{itemize}
        \item Start with prime 2 (or any prime)
        \item List all constraints on exponents
        \item Deduce values systematically
    \end{itemize}
    \item \textbf{Make It Easier}: 
    \begin{itemize}
        \item Separate problem by prime
        \item Each prime gives independent subproblem
        \item Combine results at end
    \end{itemize}
    \item \textbf{Penultimate Step}: Once exponents found, take minimum for GCD
\end{itemize}

\subsection*{C. Late-Band Strategic Playbook}
\begin{itemize}[leftmargin=*]
    \item \textbf{Answer-Choice Leverage}: 
    \begin{itemize}
        \item Test if each prime appears in GCD
        \item Only (C) has just prime 3
    \end{itemize}
    \item \textbf{Pattern Recognition}: LCM = max, GCD = min of exponents
\end{itemize}

\subsection*{D. Argument Construction Strategy}
\begin{itemize}[leftmargin=*]
    \item \textbf{Exponent Analysis Approach}: 
    \begin{enumerate}
        \item For each prime $p \in \{2, 3, 5\}$:
        \begin{itemize}
            \item Extract max constraints from six LCM equations
            \item Extract sum constraint from product equation
            \item Solve for $\nu_p(a), \nu_p(b), \nu_p(c), \nu_p(d)$
            \item Find $\min\{\nu_p(a), \nu_p(b), \nu_p(c), \nu_p(d)\}$
        \end{itemize}
        \item Combine: $\gcd(a,b,c,d) = 2^{\min_2} \cdot 3^{\min_3} \cdot 5^{\min_5}$
    \end{enumerate}
\end{itemize}

\subsection*{E. Cross-Domain Translation}
\begin{itemize}[leftmargin=*]
    \item \textbf{Number Theory to Algebra}: Exponents satisfy algebraic constraints
    \item \textbf{Constraint Satisfaction}: Logical deduction from system
\end{itemize}
\end{strategicbox}

\begin{tacticalbox}
\subsection*{A. Core Mathematical Tactics}
\begin{itemize}[leftmargin=*]
    \item \textbf{Exponent Notation}: $\nu_p(x)$ simplifies GCD/LCM analysis
    \item \textbf{Separation by Prime}: Solve independently for each prime
    \item \textbf{Constraint Analysis}: Use max and sum constraints
    \item \textbf{Logical Deduction}: Determine unique values
\end{itemize}

\subsection*{B. Domain-Specific Tactics}

\textbf{Number Theory:}
\begin{itemize}[leftmargin=*]
    \item \textbf{GCD Property}: $\gcd(x,y) = \prod p^{\min\{\nu_p(x), \nu_p(y)\}}$
    \item \textbf{LCM Property}: $\text{lcm}(x,y) = \prod p^{\max\{\nu_p(x), \nu_p(y)\}}$
    \item \textbf{Product Property}: $xy = \gcd(x,y) \cdot \text{lcm}(x,y)$
    \item \textbf{Prime Factorization}: Unique representation
\end{itemize}

\textbf{Algebra:}
\begin{itemize}[leftmargin=*]
    \item \textbf{Max/Min Equations}: Solve systems with max constraints
    \item \textbf{Linear Constraint}: Sum of exponents from product
\end{itemize}

\subsection*{C. Crossover Tactics}
\begin{itemize}[leftmargin=*]
    \item \textbf{Number Theory-Algebraic}: Exponents as variables
    \item \textbf{Logical-Mathematical}: Constraint satisfaction
\end{itemize}
\end{tacticalbox}

\begin{techniquesbox}
\subsection*{A. High-Probability Techniques}
\begin{itemize}[leftmargin=*]
    \item \textbf{Primary Technique}: Exponent Analysis by Prime
    \item \textbf{Recognition Cues}: 
    \begin{itemize}
        \item All numbers in prime factorization form
        \item Multiple LCM constraints given
        \item Only three primes involved
    \end{itemize}
    \item \textbf{Application - Complete Solution}: 
    
    \textbf{Step 1: Analyze Prime 2}
    
    Let $\nu_2(x)$ denote the exponent of 2 in $x$.
    
    From the LCM constraints:
    \begin{itemize}
        \item $\text{lcm}(b,c) = 2^1 \cdot \ldots$ implies $\max\{\nu_2(b), \nu_2(c)\} = 1$
        \item $\text{lcm}(a,b) = 2^3 \cdot \ldots$ implies $\max\{\nu_2(a), \nu_2(b)\} = 3$
        \item $\text{lcm}(b,d) = 2^2 \cdot \ldots$ implies $\max\{\nu_2(b), \nu_2(d)\} = 2$
        \item $\text{lcm}(c,d) = 2^2 \cdot \ldots$ implies $\max\{\nu_2(c), \nu_2(d)\} = 2$
        \item $\text{lcm}(a,c) = 2^3 \cdot \ldots$ implies $\max\{\nu_2(a), \nu_2(c)\} = 3$
        \item $\text{lcm}(a,d) = 2^3 \cdot \ldots$ implies $\max\{\nu_2(a), \nu_2(d)\} = 3$
    \end{itemize}
    
    From the product:
    \begin{itemize}
        \item $\nu_2(a) + \nu_2(b) + \nu_2(c) + \nu_2(d) = 6$
    \end{itemize}
    
    Analysis:
    \begin{itemize}
        \item From $\max\{\nu_2(b), \nu_2(c)\} = 1$: both $\nu_2(b), \nu_2(c) \leq 1$
        \item From $\max\{\nu_2(a), \nu_2(b)\} = 3$: since $\nu_2(b) \leq 1$, we need $\nu_2(a) = 3$
        \item From $\max\{\nu_2(b), \nu_2(d)\} = 2$: since $\nu_2(b) \leq 1$, we need $\nu_2(d) = 2$
        \item Sum: $3 + \nu_2(b) + \nu_2(c) + 2 = 6$, so $\nu_2(b) + \nu_2(c) = 1$
        \item Combined with $\max\{\nu_2(b), \nu_2(c)\} = 1$: one is 1, the other is 0
        \item So $(\nu_2(b), \nu_2(c)) \in \{(0,1), (1,0)\}$
    \end{itemize}
    
    Minimum exponent of 2: $\min\{3, 0 \text{ or } 1, 1 \text{ or } 0, 2\} = 0$
    
    \textbf{Step 2: Analyze Prime 3}
    
    Let $\nu_3(x)$ denote the exponent of 3 in $x$.
    
    From the LCM constraints:
    \begin{itemize}
        \item $\text{lcm}(a,b) = \ldots \cdot 3^2 \cdot \ldots$ implies $\max\{\nu_3(a), \nu_3(b)\} = 2$
        \item $\text{lcm}(a,c) = \ldots \cdot 3^3 \cdot \ldots$ implies $\max\{\nu_3(a), \nu_3(c)\} = 3$
        \item $\text{lcm}(a,d) = \ldots \cdot 3^3 \cdot \ldots$ implies $\max\{\nu_3(a), \nu_3(d)\} = 3$
        \item $\text{lcm}(b,c) = \ldots \cdot 3^3 \cdot \ldots$ implies $\max\{\nu_3(b), \nu_3(c)\} = 3$
        \item $\text{lcm}(b,d) = \ldots \cdot 3^3 \cdot \ldots$ implies $\max\{\nu_3(b), \nu_3(d)\} = 3$
        \item $\text{lcm}(c,d) = \ldots \cdot 3^3 \cdot \ldots$ implies $\max\{\nu_3(c), \nu_3(d)\} = 3$
    \end{itemize}
    
    From the product:
    \begin{itemize}
        \item $\nu_3(a) + \nu_3(b) + \nu_3(c) + \nu_3(d) = 9$
    \end{itemize}
    
    Analysis:
    \begin{itemize}
        \item From $\max\{\nu_3(a), \nu_3(b)\} = 2$: both $\nu_3(a), \nu_3(b) \leq 2$
        \item From $\max\{\nu_3(a), \nu_3(c)\} = 3$: since $\nu_3(a) \leq 2$, we need $\nu_3(c) = 3$
        \item From $\max\{\nu_3(a), \nu_3(d)\} = 3$: since $\nu_3(a) \leq 2$, we need $\nu_3(d) = 3$
        \item From $\max\{\nu_3(c), \nu_3(d)\} = 3$: consistent with $\nu_3(c) = \nu_3(d) = 3$
        \item Sum: $\nu_3(a) + \nu_3(b) + 3 + 3 = 9$, so $\nu_3(a) + \nu_3(b) = 3$
        \item Combined with $\max\{\nu_3(a), \nu_3(b)\} = 2$ and both $\leq 2$:
        \item Possible: $(1,2)$ or $(2,1)$
        \item Both scenarios give the same minimum
    \end{itemize}
    
    Minimum exponent of 3: $\min\{1 \text{ or } 2, 2 \text{ or } 1, 3, 3\} = 1$
    
    \textbf{Step 3: Analyze Prime 5}
    
    Let $\nu_5(x)$ denote the exponent of 5 in $x$.
    
    From the LCM constraints:
    \begin{itemize}
        \item $\text{lcm}(a,b) = \ldots \cdot 5^3$ implies $\max\{\nu_5(a), \nu_5(b)\} = 3$
        \item $\text{lcm}(a,c) = \ldots \cdot 5^3$ implies $\max\{\nu_5(a), \nu_5(c)\} = 3$
        \item $\text{lcm}(a,d) = \ldots \cdot 5^3$ implies $\max\{\nu_5(a), \nu_5(d)\} = 3$
        \item $\text{lcm}(b,c) = \ldots \cdot 5^2$ implies $\max\{\nu_5(b), \nu_5(c)\} = 2$
        \item $\text{lcm}(b,d) = \ldots \cdot 5^2$ implies $\max\{\nu_5(b), \nu_5(d)\} = 2$
        \item $\text{lcm}(c,d) = \ldots \cdot 5^2$ implies $\max\{\nu_5(c), \nu_5(d)\} = 2$
    \end{itemize}
    
    From the product:
    \begin{itemize}
        \item $\nu_5(a) + \nu_5(b) + \nu_5(c) + \nu_5(d) = 7$
    \end{itemize}
    
    Analysis:
    \begin{itemize}
        \item From $\max\{\nu_5(b), \nu_5(c)\} = 2$, $\max\{\nu_5(b), \nu_5(d)\} = 2$, $\max\{\nu_5(c), \nu_5(d)\} = 2$:
        \item All of $\nu_5(b), \nu_5(c), \nu_5(d) \leq 2$
        \item From $\max\{\nu_5(a), \nu_5(b)\} = 3$: since $\nu_5(b) \leq 2$, we need $\nu_5(a) = 3$
        \item Sum: $3 + \nu_5(b) + \nu_5(c) + \nu_5(d) = 7$
        \item So: $\nu_5(b) + \nu_5(c) + \nu_5(d) = 4$
        \item With all three $\leq 2$ and sum = 4, and all max pairs = 2:
        \item Must have two values equal to 2 and one equal to 0
        \item So one of $b, c, d$ has no factor of 5
    \end{itemize}
    
    Minimum exponent of 5: $\min\{3, \text{one is 0}\} = 0$
    
    \textbf{Step 4: Compute GCD}
    
    \begin{align*}
    \gcd(a,b,c,d) &= 2^{\min\{\nu_2\}} \cdot 3^{\min\{\nu_3\}} \cdot 5^{\min\{\nu_5\}}\\
    &= 2^0 \cdot 3^1 \cdot 5^0\\
    &= 3
    \end{align*}
\end{itemize}

\subsection*{B. Alternative Techniques}
\begin{itemize}[leftmargin=*]
    \item \textbf{Testing Answer Choices}:
    \begin{itemize}
        \item If $\gcd = 30 = 2 \cdot 3 \cdot 5$: all four numbers divisible by 30
        \item Check if consistent with product: $abcd = 30^4 \cdot (\text{other}) = 810000 \cdot (\ldots)$
        \item But $2^6 \cdot 3^9 \cdot 5^7$ has $\nu_2 = 6$, and $30^4 = 2^4 \cdot \ldots$ leaves $\nu_2 = 2$
        \item Would need remaining factor with $\nu_2 = 2$, distributed among 4 numbers
        \item But LCM constraints show one of them has $\nu_2 = 0$—contradiction!
        \item Similarly rule out other choices
    \end{itemize}
\end{itemize}

\subsection*{C. Technique Selection Justification}
\begin{itemize}[leftmargin=*]
    \item \textbf{Why Exponent Analysis}: 
    \begin{itemize}
        \item GCD/LCM naturally expressed via exponents
        \item Separates into independent subproblems
        \item Systematic and complete
    \end{itemize}
    \item \textbf{Time Efficiency}: 7-8 minutes with organized work
\end{itemize}
\end{techniquesbox}

\begin{verificationbox}
\subsection*{A. Verification Level Selection}
\begin{itemize}[leftmargin=*]
    \item \textbf{Chosen Level}: Level 3 - Targeted Spot Check (15-30 seconds)
    \item \textbf{Justification}: 
    \begin{itemize}
        \item Late-band problem (P24)
        \item Logic is systematic once understood
        \item Can verify key deductions
        \item Target confidence: 85-90\%
    \end{itemize}
    \item \textbf{Time Budget}: 20-30 seconds for verification
\end{itemize}

\subsection*{B. Verification Checklist}
\begin{itemize}[leftmargin=*]
    \item \textbf{Verify prime 2 analysis}:
    \begin{itemize}
        \item $\nu_2(a) = 3$, $\nu_2(d) = 2$, one of $b,c$ has 0 \checkmark
        \item Sum: $3 + 2 + 1 + 0 = 6$ \checkmark
        \item Min: 0 \checkmark
    \end{itemize}
    \item \textbf{Verify prime 3 analysis}:
    \begin{itemize}
        \item $\nu_3(c) = 3$, $\nu_3(d) = 3$ \checkmark
        \item $\nu_3(a) + \nu_3(b) = 3$ with both $\leq 2$ \checkmark
        \item Sum: $3 + 3 + 3 = 9$ \checkmark
        \item Min: 1 \checkmark
    \end{itemize}
    \item \textbf{Verify prime 5 analysis}:
    \begin{itemize}
        \item $\nu_5(a) = 3$ \checkmark
        \item Remaining sum to 4, with max pairs = 2 \checkmark
        \item One must be 0 \checkmark
        \item Min: 0 \checkmark
    \end{itemize}
    \item \textbf{Check answer}: $\gcd = 2^0 \cdot 3^1 \cdot 5^0 = 3$ \checkmark
\end{itemize}

\subsection*{C. Alternative Verification Methods}
\begin{itemize}[leftmargin=*]
    \item \textbf{Method 1}: Check consistency with one LCM
    \begin{itemize}
        \item If $\gcd(a,b,c,d) = 3$, then $3 | a$, $3 | b$, $3 | c$, $3 | d$
        \item This means $\min\{\nu_3(a), \nu_3(b)\} \geq 1$
        \item So $\text{lcm}(a,b)$ has $\nu_3 \geq 2$ (since max$\{\nu_3(a), \nu_3(b)\} = 2$)
        \item Given: $\text{lcm}(a,b) = 2^3 \cdot 3^2 \cdot 5^3$ has $\nu_3 = 2$ \checkmark
    \end{itemize}
    \item \textbf{Method 2}: Rule out other choices
    \begin{itemize}
        \item (A) 30: would require $\min\{\nu_2\} \geq 1$, but we found one has 0
        \item (B) 45 = $3^2 \cdot 5$: would require $\min\{\nu_3\} \geq 2$ and $\min\{\nu_5\} \geq 1$
        \item But $\min\{\nu_3\} = 1$ and $\min\{\nu_5\} = 0$
        \item (D) 15 = $3 \cdot 5$: would require $\min\{\nu_5\} \geq 1$, but we found 0
        \item (E) 6 = $2 \cdot 3$: would require $\min\{\nu_2\} \geq 1$, but we found 0
    \end{itemize}
\end{itemize}
\end{verificationbox}

\begin{solutionbox}
\subsection*{A. Step-by-Step Solution}
\begin{enumerate}[leftmargin=*]
    \item \textbf{Set Up Exponent Notation}: 
    \begin{itemize}
        \item Let $\nu_p(x)$ = exponent of prime $p$ in $x$
        \item All numbers expressed as $2^{\nu_2} \cdot 3^{\nu_3} \cdot 5^{\nu_5}$
        \item $\text{lcm}$ has $\nu_p = \max$ of exponents
        \item Product has $\nu_p = $ sum of exponents
        \item $\gcd$ has $\nu_p = \min$ of exponents
    \end{itemize}
    
    \item \textbf{Analyze Exponents of 2}: 
    \begin{itemize}
        \item From $\text{lcm}(b,c) = 2^1 \cdot \ldots$: $\max\{\nu_2(b), \nu_2(c)\} = 1$
        \item From $\text{lcm}(a,b) = 2^3 \cdot \ldots$: $\max\{\nu_2(a), \nu_2(b)\} = 3$
        \item From $\text{lcm}(b,d) = 2^2 \cdot \ldots$: $\max\{\nu_2(b), \nu_2(d)\} = 2$
        \item From product: $\nu_2(a) + \nu_2(b) + \nu_2(c) + \nu_2(d) = 6$
        \item Deduce: $\nu_2(a) = 3$, $\nu_2(d) = 2$, $(\nu_2(b), \nu_2(c)) \in \{(0,1), (1,0)\}$
        \item Therefore: $\min\{\nu_2\} = 0$
    \end{itemize}
    
    \item \textbf{Analyze Exponents of 3}: 
    \begin{itemize}
        \item From $\text{lcm}(a,b) = \ldots \cdot 3^2 \cdot \ldots$: $\max\{\nu_3(a), \nu_3(b)\} = 2$
        \item From $\text{lcm}(a,c) = \ldots \cdot 3^3 \cdot \ldots$: $\max\{\nu_3(a), \nu_3(c)\} = 3$
        \item From $\text{lcm}(a,d) = \ldots \cdot 3^3 \cdot \ldots$: $\max\{\nu_3(a), \nu_3(d)\} = 3$
        \item From product: $\nu_3(a) + \nu_3(b) + \nu_3(c) + \nu_3(d) = 9$
        \item Deduce: $\nu_3(c) = 3$, $\nu_3(d) = 3$, $\nu_3(a) + \nu_3(b) = 3$ with both $\leq 2$
        \item Possible: $(\nu_3(a), \nu_3(b)) \in \{(1,2), (2,1)\}$
        \item Therefore: $\min\{\nu_3\} = 1$
    \end{itemize}
    
    \item \textbf{Analyze Exponents of 5}: 
    \begin{itemize}
        \item From $\text{lcm}(a,b) = \ldots \cdot 5^3$: $\max\{\nu_5(a), \nu_5(b)\} = 3$
        \item From $\text{lcm}(b,c) = \ldots \cdot 5^2$: $\max\{\nu_5(b), \nu_5(c)\} = 2$
        \item All LCMs involving $b,c,d$ pairs have $\nu_5 = 2$
        \item From product: $\nu_5(a) + \nu_5(b) + \nu_5(c) + \nu_5(d) = 7$
        \item Deduce: $\nu_5(a) = 3$, and $\nu_5(b) + \nu_5(c) + \nu_5(d) = 4$
        \item With all three $\leq 2$ and sum = 4: two equal 2, one equals 0
        \item Therefore: $\min\{\nu_5\} = 0$
    \end{itemize}
    
    \item \textbf{Compute GCD}: 
    \begin{align*}
    \gcd(a,b,c,d) &= 2^{\min\{\nu_2\}} \cdot 3^{\min\{\nu_3\}} \cdot 5^{\min\{\nu_5\}}\\
    &= 2^0 \cdot 3^1 \cdot 5^0\\
    &= 3
    \end{align*}
\end{enumerate}

\subsection*{B. Alternative Approaches}

\textbf{Method 2: Elimination by Testing}

Test each answer choice:
\begin{itemize}
    \item (A) 30 = $2 \cdot 3 \cdot 5$: requires $\min\{\nu_2\} \geq 1$, but analysis shows one number has $\nu_2 = 0$
    \item (B) 45 = $3^2 \cdot 5$: requires $\min\{\nu_3\} \geq 2$, but we have $\min\{\nu_3\} = 1$
    \item (C) 3: consistent with our analysis \checkmark
    \item (D) 15 = $3 \cdot 5$: requires $\min\{\nu_5\} \geq 1$, but one number has $\nu_5 = 0$
    \item (E) 6 = $2 \cdot 3$: requires $\min\{\nu_2\} \geq 1$, but one number has $\nu_2 = 0$
\end{itemize}

\end{solutionbox}

\begin{competitionbox}
\subsection*{A. Time Management}
\begin{itemize}[leftmargin=*]
    \item \textbf{Estimated Time}: 7-10 minutes total
    \begin{itemize}
        \item Problem understanding: 1 minute
        \item Set up exponent notation: 30 seconds
        \item Analyze prime 2: 2-2.5 minutes
        \item Analyze prime 3: 2-2.5 minutes
        \item Analyze prime 5: 2-2.5 minutes
        \item Compute GCD: 30 seconds
        \item Verification: 30 seconds
    \end{itemize}
    \item \textbf{Time Checkpoints}: 
    \begin{itemize}
        \item At 2 minutes: Should understand approach
        \item At 5 minutes: Should have analyzed 2 primes
        \item At 8 minutes: Should have answer
    \end{itemize}
    \item \textbf{Priority Level}: Medium (skip if time-pressured)
    \begin{itemize}
        \item P24 is extremely difficult
        \item Requires systematic organization
        \item Only attempt if comfortable with number theory
        \item Skip if behind on time
    \end{itemize}
\end{itemize}

\subsection*{B. Risk Assessment}
\begin{itemize}[leftmargin=*]
    \item \textbf{Confidence Level}: 90\% after full analysis
    \begin{itemize}
        \item Systematic approach is sound
        \item Each prime analyzed independently
        \item Results consistent with all constraints
        \item Answer unique
    \end{itemize}
    \item \textbf{Abandonment Criteria}: 
    \begin{itemize}
        \item If after 3-4 minutes still confused, skip
        \item If analysis becomes too messy, move on
        \item Return at end if time permits
    \end{itemize}
    \item \textbf{Flag for Review}: 
    \begin{itemize}
        \item Yes, if uncertain about any deduction
        \item No, if all three primes analyzed correctly
    \end{itemize}
\end{itemize}

\subsection*{C. Answer Format}
\begin{itemize}[leftmargin=*]
    \item Multiple choice: Select (C) 3
    \item Double-check: $\min\{\nu_2\} = 0$, $\min\{\nu_3\} = 1$, $\min\{\nu_5\} = 0$ \checkmark
    \item Bubble carefully on answer sheet
\end{itemize}
\end{competitionbox}

\begin{keyinsightbox}
\textbf{Key Insight}: This problem becomes tractable by using exponent notation $\nu_p(x)$ for each prime $p$. The LCM constraints give max equations, the product constraint gives a sum equation, and the GCD requires finding minimum exponents.

Analyzing each prime separately:
\begin{itemize}
    \item \textbf{Prime 2}: Constraints force $\min\{\nu_2\} = 0$ (one number has no factor of 2)
    \item \textbf{Prime 3}: Constraints force $\min\{\nu_3\} = 1$ (all four divisible by 3)
    \item \textbf{Prime 5}: Constraints force $\min\{\nu_5\} = 0$ (one number has no factor of 5)
\end{itemize}

Therefore $\gcd(a,b,c,d) = 2^0 \cdot 3^1 \cdot 5^0 = 3$.

The beautiful structure: only prime 3 appears in all four numbers!
\end{keyinsightbox}

\begin{warningbox}
\textbf{Warning}: Common pitfalls:
\begin{enumerate}
    \item \textbf{Confusing LCM and GCD}: $\text{lcm}$ uses \textit{max} of exponents, $\gcd$ uses \textit{min}
    \item \textbf{Mixing primes}: Each prime must be analyzed independently
    \item \textbf{Arithmetic errors}: Carefully check that exponents sum correctly
    \item \textbf{Missing constraints}: Use \textit{all} six LCM equations
    \item \textbf{Logical errors}: From $\max\{x,y\} = k$, can only deduce both $\leq k$, not equality
    \item \textbf{Incomplete analysis}: Must analyze all three primes to get complete answer
\end{enumerate}
\end{warningbox}

\begin{tipbox}
\textbf{Pro Tip}: For GCD/LCM problems with prime factorizations:

\begin{enumerate}
    \item \textbf{Use exponent notation}: Write $\nu_p(x)$ for clarity
    \item \textbf{Know the formulas}:
    \begin{itemize}
        \item $\nu_p(\text{lcm}(x,y)) = \max\{\nu_p(x), \nu_p(y)\}$
        \item $\nu_p(\gcd(x,y)) = \min\{\nu_p(x), \nu_p(y)\}$
        \item $\nu_p(xy) = \nu_p(x) + \nu_p(y)$
    \end{itemize}
    \item \textbf{Separate by prime}: Each prime is independent
    \item \textbf{Organize your work}: Make a table of constraints for each prime
    \item \textbf{Use sum constraint last}: After using max constraints
    \item \textbf{Check consistency}: Sum should equal given product exponent
\end{enumerate}

For this problem:
\begin{itemize}
    \item The key is recognizing that only prime 3 has $\min \geq 1$
    \item Primes 2 and 5 both have at least one number not divisible by them
    \item The systematic exponent analysis makes this clear
    \item Organizing work by prime prevents confusion
\end{itemize}

Competition strategy: This is a sophisticated P24 that rewards systematic organization. If you're comfortable with exponent notation and GCD/LCM properties, it's very doable. Otherwise, skip and return if time permits.
\end{tipbox}

\section*{Final Answer}

By analyzing exponents of each prime separately:

\textbf{Prime 2}: $\min\{\nu_2(a), \nu_2(b), \nu_2(c), \nu_2(d)\} = 0$

\textbf{Prime 3}: $\min\{\nu_3(a), \nu_3(b), \nu_3(c), \nu_3(d)\} = 1$

\textbf{Prime 5}: $\min\{\nu_5(a), \nu_5(b), \nu_5(c), \nu_5(d)\} = 0$

Therefore:
$$\gcd(a,b,c,d) = 2^0 \cdot 3^1 \cdot 5^0 = 3$$

$$\boxed{\textbf{(C) } 3}$$

\section*{Confidence Assessment}
\begin{itemize}[leftmargin=*]
    \item \textbf{Confidence Level}: 95\% 
    \begin{itemize}
        \item Systematic exponent analysis is rigorous
        \item Each prime analyzed correctly with all constraints
        \item Results consistent: sums equal product exponents
        \item Minimum exponents correctly identified
        \item Answer unique and matches choice (C)
        \item Verification confirms no other choice works
    \end{itemize}
    \item \textbf{Verification Applied}: Level 3 - Targeted Spot Check
    \begin{itemize}
        \item Primary method: Exponent analysis by prime
        \item Verified: Each prime's constraints and deductions
        \item Cross-check: Ruled out other answer choices
        \item Consistency: All sums match product
    \end{itemize}
    \item \textbf{Flag for Review}: No
    \begin{itemize}
        \item High confidence after systematic analysis
        \item All constraints satisfied
        \item Logic is sound and complete
    \end{itemize}
    \item \textbf{Time Spent}: 
    \begin{itemize}
        \item Estimated: 7-8 minutes for complete solution
        \item Appropriate for P24 with systematic approach
        \item Organization by prime saves time
    \end{itemize}
\end{itemize}

\end{document}

